\documentclass[10pt,a4paper]{amsart}
\usepackage[margin=19mm]{geometry}
\usepackage{amsmath,amssymb,mathtools}
\usepackage[scr=rsfs,cal=euler]{mathalfa}
\usepackage{xfrac}
\usepackage[unicode,hidelinks,bookmarks=false]{hyperref}
\newcommand{\ie}{i.e.\ }
\newcommand{\cf}{cf.\ }
\newcommand{\resp}{respectively\ }
\newcommand{\Subsection}{\subsection}
\providecommand{\nobreakdash}{\nobreak\hbox{-}}
\setlength{\emergencystretch}{2em}
\multlinegap0pt
\usepackage{amssymb}
\usepackage{mathtools}
\newtheorem{lemma}{Lemma}
\newcommand{\C}{\mathbb{C}}
\renewcommand{\Im}{\mathrm{Im}\,}
\renewcommand{\Re}{\mathrm{Re}\,}
\title{Mod-$\phi$ convergence for random variables with cyclotomic generating functions}
\author{Christoph Th\"ale, Philipp Tuchel}
\date{Source: arXiv:2608.14143v1 (CC BY 4.0)}
\begin{document}
\maketitle

\noindent\emph{Source and licence.} Mod-$\phi$ convergence for random variables with cyclotomic generating functions. Version arXiv:2608.14143v1 (CC BY 4.0), distributed by arXiv under the Creative Commons Attribution 4.0 International licence, http://creativecommons.org/licenses/by/4.0/, as recorded in the arXiv licence metadata for this exact version and shown on the abstract page https://arxiv.org/abs/2608.14143v1. Full licence notice: \texttt{licenses/arxiv-2608.14143-CC-BY-4.0.txt}.\par\smallskip

\noindent\textit{Editorial scope.} Counted unit: the article's lemma Df (source0.tex lines 496--499). The article's complete original proof of it is delivered with it (source0.tex lines 500--530, one \texttt{proof} environment of 31 lines). No mathematics is written, repaired or re-derived: the statement and the proof are the article's own lines, in the article's own notation, and no retained line is substituted or re-flowed.  No further source-local context is needed: every cross-reference of every retained line resolves inside the two delivered spans.  Nothing was omitted from any retained span, so the package carries no omission record.  No retained line contains a citation, so the wrapper carries no bibliography and no bibliography entry.  No retained line contains a figure environment, an image-inclusion command or drawing code, so no image is delivered and the package declares no asset dependency.  Editorial formatting only, in addition: an AMS \texttt{amsart} preamble that restates the article's own declarations and macros used by the retained lines, namely the environments \texttt{lemma} on separate counters, together with the standard packages those lines need.  The retained environments print in this document as Lemma 1 in order of appearance, whereas the article prints the same items with the numbering of its own full document.  The complete native source of the article as distributed in the arXiv e-print for this exact version is bundled unchanged as \texttt{sources/207636/source0.tex}.
\begin{lemma}[Properties of the domain]\label{lem: properties of Df}
Let $f \in C[0, 1]$ be a non-negative function. Then the set $\mathscr{D}$ is open.
If, in addition, $f$ is monotonically increasing and $f(1)>0$, then $\widetilde{\mathscr{D}}$ is open and connected. If, furthermore, $f(0) = 0$, then $\mathscr{D}$ is star-shaped and hence connected.
\end{lemma}

\begin{proof}
Let $g(w) := \frac{1-e^{w}}{-w}$ for $w \ne 0$ and $g(0) := 1$. This function is entire. Recall that the set $\mathscr{D}$ is defined as
\[
\mathscr{D} = \{z \in \mathbb{C} \mid g(zf(t)) \notin (-\infty, 0] \text{ for all } t \in [0, 1] \text{ with } f(t)>0\}.
\]
To see that $\mathscr{D}$ is open, fix $z_0\in \mathscr{D}$ and consider the continuous map $H(z,t):=g(zf(t))$ on $\C\times[0,1]$. Because $g(z_0f(t))$ avoids the closed set $(-\infty,0]$ for every $t$, the compactness of $[0,1]$ implies that the distance
\[
\delta:=\inf_{t\in[0,1]}\,d\bigl(g(z_0f(t)),(-\infty,0]\bigr)
\]
is strictly positive. The uniform continuity of $H$ on some compact neighbourhood of $(z_0,[0,1])$ provides an $\varepsilon>0$ such that $|z-z_0|<\varepsilon$ implies $|H(z,t)-H(z_0,t)|<\delta/2$ for all $t$ and $z$ in that neighbourhood. Consequently, $g(zf(t))$ remains outside $(-\infty,0]$. Therefore $B(z_0,\varepsilon)\subset \mathscr{D}$, which shows that $\mathscr{D}$ is open.

Next, we establish connectedness, considering first the case where $f(0)$ is not necessarily zero. Assume that $f$ is monotonically increasing and set $M:=f(1)>0$. Clearly, if $\Re z<0$ then $g\bigl(zf(t)\bigr)\notin(-\infty,0]$ for all $t\in[0,1]$, so $\{z\in\C:\Re z<0\}\subset \mathscr{D}$. It remains to show that $g(w)\notin(-\infty,0]$ whenever $|\,\Im w\,|<\pi$. Using the representation $g(w)=\int_{0}^{1}e^{sw}\,ds$, we compute for $w=x+iy$,
\[
\Im g(x+iy)=\int_{0}^{1}e^{sx}\sin(sy)\,ds.
\]
Since $e^{sx}$ is positive, the sign of the integrand depends on the sign of $\sin(sy)$. If $y\in (0,\pi)$, then $sy\in (0, \pi)$ for $s\in(0,1)$, and thus $\sin(sy)$ is positive, making the integral strictly positive. Similarly, if $y\in (-\pi, 0)$, the integral is strictly negative. Hence, $g(w)\notin(-\infty,0]$ whenever $|\,\Im w\,|<\pi$. Since $0\le f(t)\le M$ for $t\in[0,1]$, it follows that $|\,\Im z\,|<\pi/M$ implies $|\,\Im(zf(t))\,|<\pi$ for all $t$, and thus $z\in\mathscr{D}$. This shows that the set
\[
\widetilde{\mathscr{D}}=\{z\in\C:\Re z<0\}\cup\{z\in\C:|\,\Im z\,|<\pi/M\}
\]
is contained in $\mathscr{D}$ and is connected.

Finally, assume $f(0) = 0$. Since $f$ is continuous and increasing, we have $f([0,1])=[0,M]$. Let $S:=\{w\in\C: g(w)\in(-\infty,0]\}$. Fix $z\in \mathscr{D}$. Then $g\bigl(zf(t)\bigr)\notin(-\infty,0]$ for all $t\in(0,1]$ with $f(t)>0$, which implies $g(zy)\notin(-\infty,0]$ for all $y\in(0,M]$. Therefore, the line segment
\[
\Gamma_z:=\{zy:0<y\le M\}
\]
is disjoint from $S$. It follows that every shorter segment
\[
\Gamma_{\lambda z}:=\{\lambda zy:0<y\le M\},\qquad 0<\lambda\le 1,
\]
is also disjoint from $S$. Hence, $\lambda z\in \mathscr{D}$ for all $\lambda\in(0,1]$, meaning the full line segment $[0,z]$ lies inside of $\mathscr{D}$. Therefore, $\mathscr{D}$ is star-shaped with respect to the origin and, in particular, connected.
\end{proof}

\end{document}
